\section{Elliptic Curve Cryptography (ECC)}

\subsection{Pengantar Elliptic Curve}

Elliptic Curve Cryptography (ECC) adalah alternatif modern untuk RSA yang menawarkan keamanan setara dengan ukuran kunci yang jauh lebih kecil \cite{hankerson2004}. ECC didasarkan pada matematika kurva eliptik atas field hingga, memberikan keamanan yang lebih efisien untuk perangkat dengan sumber daya terbatas seperti smartphone dan IoT devices.

\textbf{Keunggulan ECC vs RSA}:
\begin{itemize}
  \item \textbf{Ukuran kunci}: ECC 256-bit $\approx$ RSA 3072-bit dalam keamanan
  \item \textbf{Efisiensi komputasi}: Operasi ECC lebih cepat untuk ukuran kunci yang sama
  \item \textbf{Ukuran signature}: ECC signature lebih kecil
  \item \textbf{Implementasi}: Cocok untuk perangkat dengan sumber daya terbatas
\end{itemize}

\subsection{Kurva Eliptik atas Field Hingga}

Kurva eliptik atas field hingga $\mathbb{F}_p$ didefinisikan sebagai:
\[
E: y^2 = x^3 + ax + b \pmod{p}
\]
dimana $4a^3 + 27b^2 \neq 0 \pmod{p}$ (tidak ada singularitas).

\begin{contoh}
Kurva secp256r1 (NIST P-256):
\begin{itemize}
  \item Prime field: $p = 2^{256} - 2^{224} + 2^{192} + 2^{96} - 1$
  \item Kurva: $y^2 = x^3 - 3x + b \pmod{p}$
  \item Koefisien $b$: \path{0x5AC635D8AA3A93E7B3EBBD55769886BC651D06B0CC53B0F63BCE3C3E27D2604B}
\end{itemize}
\end{contoh}

\subsection{Struktur Aljabar Kurva Eliptik}

Himpunan titik pada kurva eliptik membentuk grup abelian dengan operasi penjumlahan titik \cite{silverman2009}:

\textbf{Elemen Identitas}: Titik di tak hingga $\mathcal{O}$

\textbf{Penjumlahan Titik}: Diberikan dua titik $P$ dan $Q$, penjumlahan $P + Q$ didefinisikan sebagai:
\begin{enumerate}
  \item Jika $P = \mathcal{O}$, maka $P + Q = Q$
  \item Jika $Q = \mathcal{O}$, maka $P + Q = P$
  \item Jika $x_P = x_Q$ dan $y_P = -y_Q$, maka $P + Q = \mathcal{O}$ (invers)
  \item Jika tidak, maka garis melalui $P$ dan $Q$ memotong kurva di titik ketiga $R$, dan $P + Q$ adalah refleksi $R$ terhadap sumbu x
\end{enumerate}

\textbf{Perkalian Titik}: Perkalian skalar $kP$ adalah penjumlahan titik $P$ dengan dirinya sendiri sebanyak $k$ kali:
\[
kP = \underbrace{P + P + \cdots + P}_{k \text{ kali}}
\]

\begin{figure}[htbp]
  \centering
  \begin{tikzpicture}[scale=0.8]
    % Draw axes
    \draw[->] (-3,0) -- (4,0) node[right] {$x$};
    \draw[->] (0,-2) -- (3,2) node[above] {$y$};
    
    % Draw elliptic curve y^2 = x^3 - 3x + 1 (non-negative for x in ~[-1.88,0.35] and [1.54,infty))
    \draw[thick, blue] plot[smooth, domain=-1.86:0.34] (\x, {sqrt(\x*\x*\x - 3*\x + 1)});
    \draw[thick, blue] plot[smooth, domain=-1.86:0.34] (\x, {-sqrt(\x*\x*\x - 3*\x + 1)});
    \draw[thick, blue] plot[smooth, domain=1.54:3.5] (\x, {sqrt(\x*\x*\x - 3*\x + 1)});
    \draw[thick, blue] plot[smooth, domain=1.54:3.5] (\x, {-sqrt(\x*\x*\x - 3*\x + 1)});
    
    % Draw points
    \node[circle, fill=red, inner sep=2pt] (P) at (1,1) {};
    \node[below right] at (P) {$P$};
    
    \node[circle, fill=red, inner sep=2pt] (Q) at (-1,1) {};
    \node[below left] at (Q) {$Q$};
    
    \node[circle, fill=green, inner sep=2pt] (R) at (2,-2) {};
    \node[below] at (R) {$R$};
    
    \node[circle, fill=green, inner sep=2pt] (PplusQ) at (2,2) {};
    \node[above] at (PplusQ) {$P+Q$};
    
    % Draw lines
    \draw[dashed, red] (P) -- (Q);
    \draw[dashed, green] (PplusQ) -- (R);
  \end{tikzpicture}
  \caption{Penjumlahan titik pada kurva eliptik: $P + Q = R'$}
\end{figure}

\subsection{Implementasi Operasi Kurva Eliptik}

\subsubsection{Representasi Titik}

Titik pada kurva eliptik direpresentasikan sebagai struktur:
\begin{lstlisting}[style=cstyle, caption={Representasi Titik ECC}]
typedef struct {
    int x;
    int y;
    int is_infinity;  // flag untuk titik di tak hingga
} ECPoint;
\end{lstlisting}

\subsubsection{Penjumlahan Titik dalam C}

\begin{lstlisting}[style=cstyle, caption={Penjumlahan Titik Kurva Eliptik}]
ECPoint point_add(ECPoint P, ECPoint Q, int p, int a, int b) {
    ECPoint R;
    
    if (P.is_infinity) return Q;
    if (Q.is_infinity) return P;
    
    // Check for invers
    if (P.x == Q.x && P.y == -Q.y % p) {
        R.is_infinity = 1;
        return R;
    }
    
    // Calculate slope
    int lambda;
    if (P.x == Q.x) {
        // Point doubling: 2P
        lambda = (3 * P.x * P.x + a) * mod_inverse(2 * P.y, p) % p;
    } else {
        // Point addition: P + Q
        lambda = (Q.y - P.y) * mod_inverse(Q.x - P.x, p) % p;
    }
    
    // Calculate result
    R.x = (lambda * lambda - P.x - Q.x) % p;
    R.y = (lambda * (P.x - R.x) - P.y) % p;
    R.is_infinity = 0;
    
    return R;
}
\end{lstlisting}

\subsubsection{Perkalian Skalar dengan Double-and-Add}

\begin{lstlisting}[style=cstyle, caption={Perkalian Skalar dengan Double-and-Add}]
ECPoint scalar_multiply(ECPoint P, int k, int p, int a, int b) {
    ECPoint result;
    result.is_infinity = 1;  // Initialize to infinity
    
    ECPoint current = P;
    
    while (k > 0) {
        if (k & 1) {
            result = point_add(result, current, p, a, b);
        }
        current = point_add(current, current, p, a, b);  // Double
        k >>= 1;
    }
    
    return result;
}
\end{lstlisting}

\subsection{Domain Parameters ECC}

Setiap kurva eliptik memiliki domain parameters yang didefinisikan secara unik:

\begin{table}[htbp]
  \centering
  \small
  \begin{tabular}{ll}
    \toprule
    Parameter & Deskripsi \\
    \midrule
    $p$ & Prime modulus field \\
    $a, b$ & Koefisien kurva $y^2 = x^3 + ax + b$ \\
    $G$ & Generator point (base point) \\
    $n$ & Order generator (jumlah titik dalam grup) \\
    $h$ & Cofactor $h = \#E(\mathbb{F}_p)/n$ \\
    \bottomrule
  \end{tabular}
  \caption{Domain parameters kurva eliptik.}
\end{table}

\begin{contoh}
Domain parameters secp256r1 (NIST P-256):
\begin{itemize}
  \item $p$ = \path{0xFFFFFFFF00000001000000000000000000000000FFFFFFFFFFFFFFFFFFFFFFFF}
  \item $a$ = \path{0xFFFFFFFF00000001000000000000000000000000FFFFFFFFFFFFFFFFFFFFFFFC}
  \item $b$ = \path{0x5AC635D8AA3A93E7B3EBBD55769886BC651D06B0CC53B0F63BCE3C3E27D2604B}
  \item $G_x$ = \path{0x6B17D1F2E12C4247F8BCE6E563A440F277037D812DEB33A0F4A13945D898C296}
  \item $G_y$ = \path{0x4FE342E2FE1A7F9B8EE7EB4A7C0F9E162BCE33576B315ECECBB6406837BF51F5}
  \item $n$ = \path{0xFFFFFFFF00000000FFFFFFFFFFFFFFFFBCE6FAADA7179E84F3B9CAC2FC632551}
  \item $h$ = 1
\end{itemize}
\end{contoh}

\subsection{ECC untuk Kriptografi Kunci Publik}

\subsubsection{ECDSA (Elliptic Curve Digital Signature Algorithm)}

ECDSA adalah adaptasi DSA untuk kurva eliptik \cite{johnson2001}:

\textbf{Key Generation}:
\begin{enumerate}
  \item Pilih bilangan acak $d$ dengan $1 < d < n$
  \item Hitung kunci publik $Q = dG$
  \item Private key: $d$, Public key: $Q$
\end{enumerate}

\textbf{Signature Generation}:
\begin{enumerate}
  \item Pilih bilangan acak $k$ dengan $1 < k < n$
  \item Hitung $R = kG = (x_R, y_R)$
  \item Hitung $r = x_R \bmod n$
  \item Hitung $s = k^{-1}(H(m) + dr) \bmod n$
  \item Signature: $(r, s)$
\end{enumerate}

\textbf{Signature Verification}:
\begin{enumerate}
  \item Hitung $w = s^{-1} \bmod n$
  \item Hitung $u_1 = H(m)w \bmod n$
  \item Hitung $u_2 = rw \bmod n$
  \item Hitung $X = u_1G + u_2Q = (x_X, y_X)$
  \item Hitung $v = x_X \bmod n$
  \item Valid jika $v = r$
\end{enumerate}

\subsubsection{ECDH (Elliptic Curve Diffie-Hellman)}

ECDH adalah protokol pertukaran kunci berdasarkan kurva eliptik \cite{menezes1996}:

\textbf{Key Exchange}:
\begin{enumerate}
  \item Alice memilih $a$ acak, hitung $A = aG$
  \item Bob memilih $b$ acak, hitung $B = bG$
  \item Alice mengirim $A$ ke Bob, Bob mengirim $B$ ke Alice
  \item Alice hitung $K = aB$, Bob hitung $K = bA$
  \item Keduanya mendapatkan kunci bersama yang sama
\end{enumerate}

\begin{lstlisting}[style=cstyle, caption={ECDH Key Exchange}]
// Alice's side
int a = random_int(1, n-1);
ECPoint A = scalar_multiply(G, a, p, curve_a, curve_b);

// Bob's side  
int b = random_int(1, n-1);
ECPoint B = scalar_multiply(G, b, p, curve_a, curve_b);

// Shared secret computation
ECPoint shared_alice = scalar_multiply(B, a, p, curve_a, curve_b);
ECPoint shared_bob = scalar_multiply(A, b, p, curve_a, curve_b);

// Both should get the same point
assert(shared_alice.x == shared_bob.x && shared_alice.y == shared_bob.y);
\end{lstlisting}

\subsection{Keamanan ECC}

Keamanan ECC bergantung pada \textbf{Elliptic Curve Discrete Logarithm Problem} (ECDLP):
\begin{itemize}
  \item Diberikan $P$ dan $Q = kP$, sangat sulit menghitung $k$
  \item Tidak ada algoritma sub-eksponensial yang diketahui untuk ECDLP
  \item Pollard's rho algorithm adalah yang tercepat dengan kompleksitas $O(\sqrt{n})$
\end{itemize}

\begin{table}[htbp]
  \centering
  \small
  \begin{tabular}{lcc}
    \toprule
    Algoritma & Ukuran Kunci & Keamanan Setara \\
    \midrule
    RSA-1024 & 1024 bit & ECC-160 \\
    RSA-2048 & 2048 bit & ECC-224 \\
    RSA-3072 & 3072 bit & ECC-256 \\
    RSA-7680 & 7680 bit & ECC-384 \\
    RSA-15360 & 15360 bit & ECC-512 \\
    \bottomrule
  \end{tabular}
  \caption{Perbandingan keamanan RSA vs ECC.}
\end{table}

\subsection{Implementasi Praktis}

\textbf{Library yang Direkomendasikan}:
\begin{itemize}
  \item OpenSSL: Implementasi lengkap ECC (ECDSA, ECDH, EdDSA)
  \item libsodium: API modern dan aman untuk kriptografi
  \item Bouncy Castle: Library Java untuk kriptografi
  \item wolfSSL: Library C yang ringan untuk embedded systems
\end{itemize}

\textbf{Pertimbangan Implementasi}:
\begin{itemize}
  \item Gunakan kurva yang telah divalidasi (secp256r1, secp384r1, curve25519)
  \item Gunakan RNG yang aman untuk kunci privat
  \item Implementasi constant-time untuk side-channel resistance
  \item Validasi input untuk mencegah invalid curve attacks
\end{itemize}

\begin{catatan}
Implementasi ECC dari contoh di atas adalah untuk tujuan edukasi. Untuk produksi, gunakan library yang telah diaudit keamanannya dan mengikuti standar industri seperti FIPS 186-4 untuk implementasi digital signature.
\end{catatan}
